\begin{theorem}[Comparison in an abelian setting, under (E) and (S) alone]
Let $\mathbb G$ be a sub-elliptic setting (\noderef{CarnotSetting}) with group law $*$,
homogeneous norm $\|\cdot\|$, step $r$, frame $\sigma$, $q:=2\,r!$ and
$d(x,y):=\|y^{-1}*x\|$, and assume the cited results of the literature
at $\mathbb G$ (\noderef{CitedResults}). Assume $\mathbb G$ is \emph{abelian}: $x*y=y*x$ for all
$x,y\in\mathbb R^{n}$. Let $A:\mathbb R^{m}\to\mathbb R^{m\times m}$ and
$H:\mathbb R^{m}\to\mathbb R$ be continuous on $\mathbb R^{m}\mns\{0\}$, with $A(\xi)$
symmetric and positive semidefinite for $\xi\ne0$, and satisfy the conditions of strict
ellipticity in the gradient direction (E) (1.9) and of homogeneity along rays (S) (1.10)
--- that is, let $(A,H)$ satisfy the structure conditions of
\noderef{StructureConditionsES}. No normalization of $(A,H)$ is assumed. Let
$\Om\subset\mathbb R^{n}$ be a bounded open set with $\mathbb R^{n}\mns\Om\ne\emptyset$, let
$u,v:\mathbb R^{n}\to\mathbb R$ be continuous on $\bar\Om$, and assume
\begin{enumerate}
\item $u$ is a viscosity sub-solution (\noderef{IsViscSubsol}) and $v$ a viscosity
  super-solution (\noderef{IsViscSupersol}) of $\mathcal L\cdot=0$ in $\Om$;
\item $u\le v$ on $\del\Om$.
\end{enumerate}
Let $R\ge0$ satisfy $|u|\le R/2$ and $|v|\le R/2$ on $\bar\Om$. Then $u\le v$ in $\Om$.
\end{theorem}
\begin{proof}
This is the conclusion the paper reaches in the first paragraph of Theorem~1.3(iii): "By the
proof of Theorem Theorem~1.2, it suffices to show that under the standing assumptions
(3.10)--(3.11) a semi-convex/semi-concave pair cannot exist." The comparison machinery is
not re-derived: the regularized pair on the reduced domain is such a pair, and
\noderef{AbelianNoStandingPair} forbids it.

Two points of the formulation first, as in Theorem~1.2.

\emph{The parameter $R$.} The paper defines $R:=2\max\{\|u\|_{L^\infty},\|v\|_{L^\infty}\}$
just before (3.15), whereas \noderef{ReducedDomain} carries $R$ as a parameter. We
therefore quantify over every admissible $R$, i.e. every $R\ge0$ with $|u|,|v|\le R/2$ on
$\bar\Om$. This contains the paper's statement: the paper's $R$ is admissible, since
$\|u\|_{L^\infty}\le\max\{\|u\|_{L^\infty},\|v\|_{L^\infty}\}=R/2$ and likewise for $v$, and
it is the smallest admissible one. It is also not stronger than the paper's statement in any
way the paper does not justify: $R$ enters only through the definition of $\Om[\eps]$, where
a larger $R$ makes $\Om[\eps]$ smaller, and the role of $|u|\le R/2$ is exactly the one the
paper gives it in the attainment step of Lemma~3.8(4), namely to make
$\|z_x\|^{q}\le4\eps\|u\|_{L^\infty}\le2R\eps<(1+4R)\eps$.

\emph{The hypothesis $\mathbb R^{n}\mns\Om\ne\emptyset$.} This is the paper's own
convention made explicit: (3.15) takes an infimum over $\mathbb R^{n}\mns\Om$, and the
paper's $\Om\subset\mathbb R^{n}$ is a proper subset. It is automatic as soon as $n\ge1$,
since no bounded subset of $\mathbb R^{n}$ is then all of $\mathbb R^{n}$.

\emph{Why the normalization is not in the statement.} The printed theorem is for general
$A,H$ satisfying (E) and (S); (2.3) is an internal step, carried out once in Step 1
below. It appears as a hypothesis of \noderef{AbelianNoStandingPair} only because
\noderef{ComparisonNondegenerateSemiconvex}, which that node uses, is stated for a
normalized pair; this is the level at which it is discharged.

\medskip
Suppose the conclusion fails, so $u(x_0)>v(x_0)$ for some $x_0\in\Om$; we derive a
contradiction.

\medskip\noindent\textbf{Step 0: $M_0>0$.} $\Om$ is bounded, so $\bar\Om$ is compact, and it
is non-empty ($x_0\in\Om\subseteq\bar\Om$); $u-v$ is continuous on $\bar\Om$, so it attains
its maximum there at some $x_M\in\bar\Om$ and
\[
 M_0:=u(x_M)-v(x_M)\ \ge\ u(x_0)-v(x_0)\ >\ 0 .
\]
Put $\tau:=M_0/2>0$.

\medskip\noindent\textbf{Step 1: normalize once.} The only hypothesis of
\noderef{NormalizationExists} is that $(A,H)$ satisfy the structure conditions, which we
have. It provides $A_0,H_0$ and $a$ with $a(\xi)>0$, $A(\xi)=a(\xi)A_0(\xi)$ and
$H(\xi)=a(\xi)H_0(\xi)$ for $\xi\ne0$, with $(A_0,H_0)$ satisfying both the structure
conditions (\noderef{StructureConditionsES}) and the normalization (2.3)
(\noderef{NormalizedAH}). Write $\mathcal L_0$ for the operator built from $(A_0,H_0)$.

The hypotheses of \noderef{NormalizedSameViscositySolutions} are exactly those three
identities, and its open set and function are arbitrary parameters; applying it with $\Om$
and $u$, and with $\Om$ and $v$, gives that $u$ is a viscosity sub-solution and $v$ a
viscosity super-solution of $\mathcal L_0\cdot=0$ in $\Om$.

\medskip\noindent\textbf{Step 2: the regularized pair.} All hypotheses of
\noderef{RegularizedPairStandingAssumptions} hold for
$(\mathbb G,\text{cited results at }\mathbb G,A_0,H_0,\Om,u,v,R,x_M)$: the structure
conditions on
$(A_0,H_0)$ by Step 1; $\Om$ open, bounded, with non-empty complement; $u,v$ continuous on
$\bar\Om$ and a viscosity sub-/super-solution pair for $\mathcal L_0$ in $\Om$ by Step 1;
$u\le v$ on $\del\Om$; $R\ge0$ with $|u|,|v|\le R/2$ on $\bar\Om$; $x_M$ a maximum point of
$u-v$ over $\bar\Om$; and $M_0>0$. Taking $\tau=M_0/2$ and
$E=\argmax_\Om(u-v)$ and writing
\[
 u^{\eps}:=\mathrm{SupConv}\ \mathbb G\ \Om\ u\ \eps\ (\noderef{SupConv}),\qquad
 \tilde v_{\eps}:=\big(y\mapsto\mathrm{InfConv}\ \mathbb G\ \Om\ v\ \eps\ (y)+\tau\big)
 \ (\noderef{InfConv}),
\]
and $\Om[\eps]$ for the reduced domain (\noderef{ReducedDomain}) with this $R$, it yields an
$\eps_1>0$ such that for every $\eps\in(0,\eps_1)$:
\begin{itemize}
\item[(b)] $u^{\eps}$ and $\tilde v_{\eps}$ are continuous on $\overline{\Om[\eps]}$;
\item[(c)] $u^{\eps}$ is locally semi-convex and $\tilde v_{\eps}$ locally semi-concave on
  $\Om[\eps]$, i.e. $-\tilde v_{\eps}$ is locally semi-convex there
  (\noderef{LocallySemiConvexOn});
\item[(d)] $u^{\eps}$ is a viscosity sub-solution and $\tilde v_{\eps}$ a viscosity
  super-solution of $\mathcal L_0\cdot=0$ in $\Om[\eps]$;
\item[(e)] $u^{\eps}-\tilde v_{\eps}\le-\tau/2$ on $\del\Om[\eps]$;
\item[(f)] there is $x_e\in\Om[\eps]$ at which $u^{\eps}-\tilde v_{\eps}$ attains its maximum
  over $\overline{\Om[\eps]}$, with $\tau\le u^{\eps}(x_e)-\tilde v_{\eps}(x_e)$.
\end{itemize}
(It also yields items 1, 2(a), 2(g) and 3, which we do not need here.) Fix
\[
 \eps:=\frac{\eps_1}2\in(0,\eps_1),
\]
and write $\Om':=\Om[\eps]$, $U:=u^{\eps}$, $V:=\tilde v_{\eps}$ and
$\tau':=\tau/2>0$.

\medskip\noindent\textbf{Step 3: the geometry of $\Om'$.} By items (1), (2) and (4) of
\noderef{ReducedDomainProperties} --- applicable since $\Om$ is open and bounded with
non-empty complement and $R\ge0$ --- the set $\Om'$ satisfies $\Om'\subseteq\Om$, is bounded,
and is open (here $\eps>0$). Hence $\overline{\Om'}$ is compact. Moreover
$\mathbb R^{n}\mns\Om'\ne\emptyset$: taking $p\notin\Om$ from our hypothesis, $p\notin\Om'$
because $\Om'\subseteq\Om$.

\medskip\noindent\textbf{Step 4: the (3.11) package, built as a step function.} We must
supply \noderef{AbelianNoStandingPair} with non-decreasing $d$-moduli $\om_U,\om_V$ for $U$
and $V$ on $\overline{\Om'}$ and a radius $\rho_0>0$ with
$\om_U(\rho_0)+\om_V(\rho_0)\le\tau'/4$. The paper asserts at (3.11) that such moduli exist
"because $\bar\Om$ is compact and $|x-y|\le c\,d(x,y)$" and does not construct them; a
supremum $\om(s)=\sup\{|w(a)-w(b)|:d(a,b)\le s\}$ would need a separate boundedness argument,
and none is needed, because a \emph{two-valued step function} already does everything the two
consumers ask of $\om$: it is non-decreasing, it dominates the oscillation at every scale,
and it is as small as we like at one prescribed radius. We construct it.

\emph{The construction.} Let $w:\mathbb R^{n}\to\mathbb R$ be continuous on the compact set
$K:=\overline{\Om'}$, let $M$ be such that $|w|\le M$ on $K$ (such an $M$ exists: $w$ is
continuous on the compact $K$), and let $\eta>0$. Then there are $\rho>0$ and a
non-decreasing $\om:\mathbb R\to\mathbb R$ with
\[
 |w(a)-w(b)|\le\om\big(d(a,b)\big)\quad(a,b\in K),\qquad \om(\rho)\le\eta . \tag{1}
\]
Indeed: $w$ is uniformly continuous on the compact $K$, so there is $\varrho>0$ with
$|w(a)-w(b)|\le\eta$ whenever $a,b\in K$ and $|a-b|<\varrho$. The set $K$ is bounded, so
$\nu:=1+\sup\{|a|:a\in K\}$ is a finite real number with $\nu>0$ and $|a|\le\nu$ for every
$a\in K$ (the supremum of the empty set being $0$, which only makes $\nu$ larger); by the
Hölder equivalence clause of
\noderef{CarnotSetting} at this $\nu$ there is $c>0$ with $c^{-1}|a-b|\le d(a,b)$ whenever
$|a|\le\nu$ and $|b|\le\nu$, i.e.
\[
 |a-b|\le c\,d(a,b)\qquad(a,b\in K). \tag{2}
\]
Put
\[
 \rho:=\frac{\varrho}{2c}>0,\qquad
 \om(s):=\begin{cases}\eta,& s\le\rho,\\ \eta+2|M|+1,& s>\rho.\end{cases}
\]
$\om$ is non-decreasing, because $2|M|+1>0$ and $\om$ takes the smaller value exactly on the
lower half-line. If $a,b\in K$ with $d(a,b)\le\rho$ then by (2)
$|a-b|\le c\rho=\varrho/2<\varrho$, so $|w(a)-w(b)|\le\eta=\om(d(a,b))$. If $d(a,b)>\rho$
then $\om(d(a,b))=\eta+2|M|+1$, while
$|w(a)-w(b)|\le|w(a)|+|w(b)|\le2M\le2|M|<\eta+2|M|+1$. Finally $\om(\rho)=\eta$. This proves
(1).

\emph{The instantiation.} Apply the construction twice on $K=\overline{\Om'}$, with
$\eta:=\tau/16>0$: to $w:=U$ (continuous on $\overline{\Om'}$ by (b)), giving $\rho_U>0$ and
$\om_U$; and to $w:=V$ (continuous on $\overline{\Om'}$ by (b)), giving $\rho_V>0$ and
$\om_V$. Put $\rho_0:=\min\{\rho_U,\rho_V\}>0$. Since $\rho_0\le\rho_U$ and $\om_U$ is
non-decreasing, $\om_U(\rho_0)\le\om_U(\rho_U)\le\tau/16$, and likewise
$\om_V(\rho_0)\le\tau/16$. Hence
\[
 \om_U(\rho_0)+\om_V(\rho_0)\ \le\ \frac\tau8\ =\ \frac{\tau/2}4\ =\ \frac{\tau'}4 ,
 \tag{3}
\]
which is (3.11) for the standing constant $\tau'=\tau/2$.

\medskip\noindent\textbf{Step 5: the standing assumptions for $(U,V)$ on $\Om'$, and the
contradiction.} We now check every hypothesis of \noderef{AbelianNoStandingPair} for the
data
$(\mathbb G,\text{cited results at }\mathbb G,A_0,H_0,\Om',U,V,\om_U,\om_V,\tau',x_e,\rho_0)$.
\begin{itemize}
\item The cited results at $\mathbb G$: our hypothesis, and $\mathbb G$ is the setting that
  node is applied to. Commutativity $x*y=y*x$: our hypothesis.
\item $(A_0,H_0)$ satisfies the structure conditions and (2.3): Step 1. This is where the
  normalization is spent.
\item $\Om'$ is open and bounded with $\mathbb R^{n}\mns\Om'\ne\emptyset$: Step 3.
\item $U,V$ are continuous on $\overline{\Om'}$: (b).
\item $U$ is locally semi-convex and $V$ locally semi-concave in $\Om'$: (c).
\item $U$ is a viscosity sub-solution and $V$ a viscosity super-solution of
  $\mathcal L_0\cdot=0$ in $\Om'$: (d).
\item $\om_U,\om_V$ are non-decreasing and are $d$-moduli for $U,V$ on $\overline{\Om'}$:
  Step 4.
\item $\tau'>0$ and $U-V\le-\tau'$ on $\del\Om'$: $\tau'=\tau/2>0$, and (e) is exactly
  $U-V\le-\tau/2=-\tau'$ on $\del\Om'$.
\item $x_e\in\overline{\Om'}$ is a point at which $U-V$ attains its maximum over
  $\overline{\Om'}$, and $\tau'\le U(x_e)-V(x_e)$: by (f), $x_e\in\Om'\subseteq\overline{\Om'}$
  is such a maximum point with $\tau\le U(x_e)-V(x_e)$, and $\tau'=\tau/2\le\tau$ since
  $\tau>0$.
\item $\rho_0>0$ and $\om_U(\rho_0)+\om_V(\rho_0)\le\tau'/4$: Step 4, (3).
\end{itemize}
\noderef{AbelianNoStandingPair} therefore applies and yields a contradiction. Since the only
assumption made beyond the hypotheses of the present statement was $u(x_0)>v(x_0)$ for some
$x_0\in\Om$, that assumption is false: $u(x)\le v(x)$ for every $x\in\Om$, i.e. $u\le v$ in
$\Om$.
\end{proof}
